% Limitations. Sources: HTML §8; refs/docs/ts_benchmark.md §6; refs/docs/plan.md §8;
% supp.json["ctx"] (context correction: 6.5% / 7.3%, 93.6%).

\section{Limitations}
\label{sec:limitations}

\paragraph{False alarms with many talkers or short sentences.}
With four other people speaking and the target absent, the false-alarm rate of AmphionASR~2.0 is 8.3\%; when a similar-sounding person says one short sentence and the target is absent, it is 22.7\%.

\paragraph{An empty output is not speaker verification.}
With a very poor enrollment or a very low SNR, a target-present sample may still yield an empty output, and a target-absent sample may still yield a false alarm.
Applications that need strict identity checks should add a speaker-verification step.

\paragraph{Regression in general recognition.}
AmphionASR~2.0 remains behind the model it was initialized from, Qwen3-ASR-1.7B, on WenetSpeech, WenetSpeech-Chuan, the KeSpeech dialects and English (\Cref{subsec:exp_general}).

\paragraph{Contextual correction.}
An in-house contextual-correction test has a speaker clarify an earlier word in the same utterance, for example by spelling out which character of a name is meant.
AmphionASR~2.0 rewrites the earlier mention according to the clarification only 6.5\% of the time (Qwen3-ASR-1.7B: 7.3\%), although 93.6\% of the clarifications themselves are transcribed correctly.
In TS-ASR mode, when the voice in the mixture is a text-to-speech clone of a person and the enrollment is that person's real recording, 11.6\% of samples produce an empty output.
% TODO: verify the 11.6% TTS-clone figure; it comes from the published HTML §8
% and was not found in supp.json["ctx"].

\paragraph{Near-homophone distractors in hotword lists.}
When the candidate list comes from lexicon retrieval and contains homophones and near-homophones, AmphionASR~2.0 writes wrong candidates into Mandarin output.
On real speech, its false-alarm rate is 21.7\% with the correct word in the list and 39.8\% without it, roughly twice that of Qwen3-ASR-1.7B in both cases (\Cref{subsec:exp_hotwords}).
In practice, the list should be kept short, with few near-homophones.

\paragraph{Punctuation and length.}
About 29\% of plain-ASR outputs carry punctuation, so post-processing is needed to punctuate the rest.
% TODO: verify the 29% punctuation share; it comes from the published HTML §8.
One request carries at most 30\,s of audio, and longer recordings must be segmented first.

\paragraph{Benchmark coverage.}
\tsbench is Mandarin only.
Its synthetic part uses read speech, without the prosody and interruptions of real conversation, and although its noise floors and room responses are partly real recordings, the mixing is synthetic.
The real part covers only meetings, phone calls and web conversations, with no in-car, outdoor or on-device recordings, and has only 100--460 samples per corpus.
References for the web conversations come from automatic diarization and recognition and contain some errors.
The training splits of the real meeting corpora, with different recordings and speakers, are used in training, so the real part measures in-domain behavior.
AISHELL-1 test text overlaps with AISHELL-2, which is used in training, and MagicData-RAMC assigns several speaker IDs to one person in some cases, so the real-call subset may contain training speakers.
